% BOOK-FILE-ID: f-0140c19b880e
% subsection: 階差等比複合型 \, $\displaystyle a_{n+1}=pa_n+f(n)$

\subsection{\texorpdfstring{階差等比複合型  $a_{n+1}=pa_n+f(n)$}{階差等比複合型}}\label{節：階差等比複合型}
\subsubsection{\texorpdfstring{多項式を引く  $a_{n+1}=pa_n+f(n)$}{多項式を引く}}
まずは観察から入ろう.
\[
  a_{n+1}=pa_n+f(n)
\]
これは特性方程式型漸化式の$q$の部分が定数ではなく関数になったものとして解釈できる.
このように,漸化式の右辺で$a_n$を含まずに足されている$q$や$f(n)$のような項を\textbf{\bookterm{addend}{付加項}}と呼ぶ.
\sectionref*{節：特性方程式型}では,不動点の値を引く手法と,隣接2項間の差を取る手法,$p$で割る手法で漸化式を等比型や階差型に帰着させた.
今回もその方針で考えていこう.

\paragraphbf{不動関数を引く}

不動関数は本書が導入する独自の用語である.

特性方程式型漸化式の解法は
\[
  a_{n+1}-\alpha=p(a_n-\alpha)
\]
と置くことから始めた.
ここでは,不動点を$\alpha$と仮置きしていた.

階差等比複合型でも同様の手法を取ってみよう.
特性方程式型との違いは定数ではなく関数$f(n)$が足されていることにあるため,同様の手法を取るためには引く値も関数でなくてはならない.
したがって
\[
  a_{n+1}-g(n+1)=p\left\{a_n-g(n)\right\}
\]
を満たすような$g(n)$を考えてやればよい.
ここで$p\ne1$なら,$g(n+1)-pg(n)$の最高次の係数は$g(n)$の最高次の係数の$1-p$倍である.
したがって,$f(n)$が
\[
  f(n)=\sum_{k=0}^{m} x_kn^k
\]
で表される時,
\[
  g(n)=\sum_{k=0}^{m} y_kn^k
\]
と置けばよい.
あとはこれを満たすような定数$y_k$を$x_k$それぞれと係数比較\footnote{このように未知の係数を比較して求める方法を未定係数法という.}することによってその解を得ることができ,等比型に帰着する.
$p=1$のときは階差型であり,$f(n)$が$m$次なら,引く多項式は$m+1$次になる.

\begin{bookpoint}
  \textbf{POINT}
  左辺を$n+1$に関する式,右辺を$n$に関する式に変形し,等比型に帰着させる.
\end{bookpoint}

\paragraphbf{階差を取る}

ここで扱う解法は使える場面が限られるが,この解法ができないのとやらないのは大きく意味が異なる.
他の手法のところで活きる発想でもあるので,きちんと身につけておこう.

そもそもこの漸化式が難しく見えるのは$f(n)$などというややこしいものが原因である.
したがって$f(n)$を消すように式変形をすれば話は早い.
  \BookMathLayout{\[
    \begin{array}{rrl}
      &a_{n+2}={}&pa_{n+1}+f(n+1)\\
      -\,)&a_{n+1}={}&pa_n+f(n)\\ \hline
      &a_{n+2}-a_{n+1}={}&p(a_{n+1}-a_n)\\
      &&\qquad+f(n+1)-f(n)
    \end{array}
  \]}{\[
    \begin{array}{rrl}
      &a_{n+2}={}&pa_{n+1}+f(n+1)\\
      -\,)&a_{n+1}={}&pa_n+f(n)\\ \hline
      &a_{n+2}-a_{n+1}&\\
      &\multicolumn{2}{l}{{}=p(a_{n+1}-a_n)}\\
      &\multicolumn{2}{l}{\qquad+f(n+1)-f(n)}
    \end{array}
  \]}
  ここで$f(n)$が1次であれば,定数項しか残らないので,$b_n=a_{n+1}-a_n$と置くことでシンプルな特性方程式型になる.
  
  $f(n)$が2次以上の場合は,その次数を$m$とすると$f(n+1)-f(n)$は$m-1$次になり,1回階差を取っただけでは付加項が残る.
  ただし,$b_n=a_{n+1}-a_n$も再び階差等比複合型であり,付加項の次数は1つ下がっている.
  したがって階差を取る操作を$m$回繰り返せば付加項は定数になり,特性方程式型に帰着する.
  その代わり,元の$a_n$へ戻すには和を$m$回取る必要があるので,実際に計算するのは$f(n)$が1次の場合に限られるだろう.
  
\paragraphbf{指数で割る}

  特に
  \[
    f(n)=rq^n
  \]
  である時にはここで扱う解法が有効である.ここでの$q$は指数の底であり,先ほど特性方程式型の定数項を表した$q$とは別である.
  以下では$q\ne0$とする.\footnote{不動関数を引いても解ける.$q\ne p$なら$\frac{rq^n}{q-p}$を,$q=p$なら$rnp^{n-1}$を引けばよい.}
  
   結果から言えば,両辺を$q^{n+1}$で割って,
   \[
      \frac{a_{n+1}}{q^{n+1}}=\frac{p}{q}\cdot\frac{a_n}{q^n}+\frac{r}{q}
   \]
   $b_n=\frac{a_n}{q^n}$と置くことで
   \[
      b_{n+1}=\frac{p}{q}b_n+\frac{r}{q}
   \]
   となるため,$p\ne q$なら特性方程式型に帰着する.
   $p=q$なら$b_{n+1}=b_n+\frac rq$となり,等差型である.
